% Central result section: the spatial-confounding decomposition (Gelbach toggle).
% Pairs with the LEACE erasure section (sec:leace-12): erasure shows the
% embedding ENCODES geography; this section shows the price effect SURVIVES
% adjusting for it. Numbers: confounded share <5% in 11/12; NYC -0.119 -> -0.152
% (~28% of its effect, no flip), pc1_geo_r2 0.42; pc1_geo_r2 0.03-0.42, >0.10 in
% six markets. Exhibit: Figure fig:decomposition (left = naive vs geo dumbbell,
% right = pc1_geo_r2). Written to the academic-prose standard.
% Place after sec:leace-12 in assembly.

\subsection{How much of the text-price effect is location}
\label{sec:decomposition}

The erasure diagnostic of Section~\ref{sec:leace-12} establishes that the
embedding carries location, and carries it abundantly, since a held-out ridge
recovers the latitude-longitude pair from the raw representation at an $R^2$ that
climbs to one half in the largest market before the projection drives it to the
floor. Taken on its own that fact invites the deflationary reading, that a
treatment built from such a representation is geography in disguise and its
apparent price effect is spatial confounding waiting to be exposed. The inference
does not actually follow, because a representation can encode a confounder along
directions that have little to do with the direction in which it predicts the
outcome, and whether the two coincide is an empirical matter rather than something
an erasure $R^2$ can settle on its own. Answering it means holding the treatment
fixed and asking what the location it shares with price is genuinely worth to the
estimate.

The instrument for that question is a fixed-treatment decomposition in the manner
of \citet{gelbach2016covariates}, whose appeal is that the change in a coefficient
when a block of controls is added does not depend on the order in which the blocks
enter, so the share it assigns to geography is well defined rather than an
artifact of sequencing. We estimate the per-standard-deviation effect of the
leading text direction on log price twice, once controlling for the structured
confounders alone and once adding a flexible location basis of latitude,
longitude, their squares, and their interaction, and we read the gap between the
two as the portion of the effect that location explains. The comparison appears in
the left panel of Figure~\ref{fig:decomposition}, where the two estimates sit
almost on top of one another in eleven of the twelve markets, the location
adjustment shifting the effect by less than five percent of its size everywhere
but one. Philadelphia, Chicago, Atlanta, and the other high-effect markets give up
essentially nothing to the geography control, the small-effect markets give up
nothing they had to begin with, and the per-market premium that the text channel
identifies comes through the adjustment very nearly untouched by whether the model
is told where each property sits.

A reader entitled to be skeptical will ask at once whether the test has any power,
since a treatment that carried no geography in the first place would pass it for
free, and the right panel of Figure~\ref{fig:decomposition} exists to meet that
objection before it hardens into one. The share of the treatment's own variance
that the location basis explains is nowhere near zero across the panel, reaching
$0.42$ in New York and exceeding a tenth in six markets, so the leading text
direction does in fact encode location through much of the country. That the price
effect survives adjustment even where the treatment is substantially geographic is
the whole of the result, because it means the geographic component of the language
is real, the model can see it, and it still does not carry the effect of the
language on price. The worry that has hung over this literature is therefore one
that can be answered rather than only raised, and the answer across eleven of
twelve markets is that the text is doing more than reading the map back to itself.

New York is where the adjustment does bite, and it bites in precisely the place
the diagnostic says it should. The market whose text direction is the most
geographic of the twelve, at an $R^2$ of $0.42$, is also the market where folding
in location moves the effect most, from $-0.119$ to $-0.152$, a shift of about a
quarter of its magnitude that leaves the sign intact, so that the place with the
most location in its language is the place that location matters most for the
estimate. We read this less as a crack in the broader finding than as its
mechanism caught in the act, and as a reminder that the protocol earns its keep by
flagging the markets where a spatial audit is most warranted rather than by
declaring that none is ever needed. The conclusion travels only as far as its
instruments reach, and they reach less far than one would like, since the location
basis is a smooth coordinate polynomial laid over roughly thirty structured
controls that already carry census-tract demography, which leaves confounding
operating at a finer spatial grain than the polynomial can bend to, or running
through pathways those controls do not capture, outside what the comparison can
see, and since the erasure that certifies the representation's geography is itself
linear, so that a nonlinear geographic signal no ridge probe would surface could
in principle persist into the residual we are calling semantic.
